% Done
\section{Specifikace funkcí aplikace}

% Review
V této sekci se budu věnovat podrobnější specifikaci funkcí navrhované aplikace. Nejprve se zaměřím na popis jednotlivých aktérů, kteří budou s aplikací interagovat. Následně na základě funkčních požadavků sepíšu seznam případů užití, které budou podrobněji specifikovat, jak mají jednotlivé části aplikace fungovat z pohledu uživatele. Nakonec na základě případů užití sepíšu podrobnou specifikaci vybraných funkcí v jazyce Gherkin. Tato specifikace umožní efektivní spouštění automatizovaných testů pomocí nástroje Cucumber a bude také sloužit jako dokumentace toho, jak se aplikace skutečně chová. Mimo jiné bude možné tuto specifikaci využít také jako zadání pro AI agenty, kteří mohou v budoucnu na základě této specifikace přesně implementovat požadované funkce, nebo například generovat dokumentaci a testy.

% Review
\subsection{Popis aktérů}

% Review
V rámci specifikace funkcí je třeba popsat tzv. aktéry, kteří reprezentují role, které mohou s aplikací interagovat a vykonávat v ní určité akce. Tito aktéři budou využiti jak v popisu případů užití, tak ve specifikaci funkcí pomocí jazyka Gherkin. Na základě požadavků aplikace jsem identifikoval následující aktéry.

% Review
\subsubsection{Uživatel}

% Review
Uživatel je osoba, která využívá aplikaci k učení se programování. Mezi jeho hlavní aktivity patří například zapisování si kurzů, učení se programování, interakce s ostatními uživateli a interakce s prvky gamifikace v aplikaci.

% Review
\subsubsection{Premium uživatel}

% Review
Premium uživatel je osoba, která má aktivní předplatné nebo zkušební verzi předplatného. Premium uživatel může využívat všechny funkce jako uživatel. Kromě těchto akcí může vykonávat i další akce, ke kterým nemá běžný uživatel přístup.

% Review
\subsubsection{Korektor}

% Review
Korektor je osoba zodpovědná za kontrolu kvality výukového obsahu, aby byl správný a srozumitelný. V rámci aplikace může používat funkce podobně jako premium uživatel. Kromě toho může využívat podrobnější nahlašování chyb při učení.

\subsection{Seznam případů užití}

% Vysvětlit, co jsou případy užití, k čemu jsou dobré a proč je vhodné je specifikovat.
% Odkázat se na nějakou odbornou literaturu.

% Popsat postupně všechny use cases případy užití ve stejném pořadí jako jsou funkční požadavky

\subsection{Specifikace funkcí pomocí jazyku Gherkin}

% Popsat co je Gherkin, Cucumber a jaké výhody přináší. Říct, že příchod umělé inteligence zásadně změnil způsob vývoje a že jsem se rozhodl použít Gherkin, který umožní definovat přesné scénáře pro AI agenty a zároveň bude sloužit pro automatizaci testů
% Popsat, co obsahuje Gherkin feature file a přidat ukázku
% Popsat stručně, že gherkin feature files mohou sloužit jako přesná specifikace pro AI agenty, kteří pomáhají při implementaci samotných funkcí. Jelikož se jedná o přesnou specifikaci.
% Popsat, stručně, jak lze gherkin feature files následně pomocí Cucumber využít pro spuštění testů. Nepopisovat jak to funguje, jen stručně říct, že Cucumber podle jednotlivých kroků scénářů následně spouští automatizované testy a zajišťuje tak, že jsou všechny ty kroky scénářů pokryty testy a že tak aplikace skutečně funguje.

% Konečně bude tato specifikace sloužit jako dokumentace toho, jak přesně se aplikace chová a jakým způsobem ji lze využívat.